\section{CONCLUSÃO}

\subsection{Síntese das Decisões de Projeto}

O presente Trabalho Final consolidou a concepção, especificação arquitetural e modelagem integral de um sistema web de gestão clínica, o \textbf{MedAgenda}, abordando todas as fases que compõem a disciplina de Engenharia de Software I \cite{sommerville2011}. O desenvolvimento do projeto permitiu comprovar na prática a aplicação dos seguintes preceitos metodológicos:
\begin{itemize}
    \item \textbf{Arquitetura Limpa e Desacoplada}: A escolha da arquitetura em camadas combinada com o padrão MVC permitiu segregar com clareza a camada de apresentação em Single Page Application (\textbf{React.js e Vite}) do núcleo transacional transacionado via API RESTful (\textbf{Node.js, Express, Prisma ORM e PostgreSQL}) \cite{martin2019, fowler2002}. Essa decisão eliminou dependências cíclicas e blindou as regras clínicas de negócio contra mudanças na tecnologia visual.
    \item \textbf{Containerização e Reprodutibilidade Total}: A adoção estratégica do \textbf{Docker e Docker Compose} como plataforma de empacotamento e orquestração dos serviços (banco de dados relacional com volume persistente, servidor da API REST com bootstrap automático e servidor Nginx de alta performance para o frontend) resolveu definitivamente o desafio da portabilidade de ambientes e da reprodutibilidade de testes na Engenharia de Software \cite{docker2023, burns2018}. A capacidade de inicializar o ecossistema complexo via comando único (\texttt{docker compose up --build}) atesta o rigor operacional do projeto.
    \item \textbf{Rastreabilidade Integral}: A especificação estruturada de doze requisitos funcionais e seis não funcionais manteve rastreabilidade bidirecional exata e auditável com o Diagrama de Casos de Uso, os doze casos de teste prioritários e a matriz transacional de verificação e validação \cite{myers2011}.
    \item \textbf{Ergonomia e Usabilidade Clínica}: A projeção dos requisitos estruturais das cinco interfaces interativas demonstrou o compromisso com a redução de carga cognitiva e a prevenção de falhas operacionais em ambiente médico, aplicando diretamente as heurísticas clássicas de Jakob Nielsen \cite{nielsen1993, norman2013}.
\end{itemize}

\subsection{Dificuldades Encontradas e Soluções Adotadas}

Durante o desenvolvimento deste trabalho, a equipe superou desafios técnicos e analíticos relevantes:
\begin{itemize}
    \item \textbf{Modelagem de Regras Temporais no Domínio Médico}: A verificação de disponibilidade horária e a aplicação de restrições transacionadoras para cancelamento (regra das 24 horas de antecedência) exigiram cuidado na especificação dos atributos relacionais em \texttt{HorarioDisponivel} e \texttt{Consulta}, garantindo a acurácia dos cálculos temporais nos diagramas de sequência \cite{larman2004}.
    \item \textbf{Orquestração de Containerização Heterogénea}: A comunicação em rede interna do Docker entre um frontend estático servido via Nginx e uma API REST em Node.js exigiu a configuração de regras de proxy reverso (\texttt{try\_files} e repasse da rota \texttt{/api}), além de condicionar a inicialização do backend à aprovação de integridade do banco de dados relacional (\texttt{depends\_on: condition: service\_healthy}) \cite{burns2018}.
    \item \textbf{Manutenção do Foco Conceitual de Engenharia de Software}: Houve rigorosa atenção para não desviar o relatório do seu objetivo central — a especificação estrutural, modelagem arquitetural e levantamento formal de requisitos — evitando a inserção de imagens ou capturas de telas prontas, conforme orientação metodológica da disciplina \cite{sommerville2011}.
\end{itemize}

\subsection{Lições Aprendidas}

A realização do projeto MedAgenda proporcionou ensinamentos valiosos para a formação técnica e acadêmica dos integrantes da equipe:
\begin{itemize}
    \item \textbf{O Valor do Desacoplamento Arquitetural}: Compreendeu-se que o uso consciente de princípios de design e padrões de projeto (tais como os padrões GoF, GRASP e Injeção de Dependência) simplifica drasticamente a evolução transacional do sistema, o reaproveitamento de código e a criação de suítes de testes automatizados \cite{gamma1995, larman2004}.
    \item \textbf{O Poder da Containerização na Engenharia de Software}: Evidenciou-se que dominar ferramentas de virtualização leve como o Docker é indispensável para garantir que o software projetado em teoria opere de forma imutável, reprodutível e robusta em ambientes reais de homologação e produção \cite{docker2023}.
    \item \textbf{Usabilidade como Requisito Crítico em Sistemas de Saúde}: A experiência consolidou a visão de que, em sistemas de missão crítica como os da área médica, a clareza da interface, o feedback visual dos estados e a prevenção proativa de erros humanos são tão vitais quanto a performance computacional e a integridade criptográfica da aplicação \cite{nielsen1993, silberschatz2020}.
\end{itemize}
